\chapter*{Abstract}
\addcontentsline{toc}{chapter}{Abstract}

This project investigates a wearable center-fed half-wave dipole operating at \(1.120~\mathrm{GHz}\) near a simplified finite skin-fat-muscle body model. The theoretical free-space half-wave length was \(133.93~\mathrm{mm}\). Full-wave tuning reduced the total dipole length to \(122.0~\mathrm{mm}\), yielding a minimum \(S_{11}\) of \(-15.52702~\mathrm{dB}\) at \(1.1168~\mathrm{GHz}\) and \(S_{11}(1.120~\mathrm{GHz})=-15.48433~\mathrm{dB}\).

Body loading degraded both impedance matching and radiating performance. At separations of \(10~\mathrm{mm}\) and \(5~\mathrm{mm}\), the dominant matching minima shifted to \(1.0312~\mathrm{GHz}\) and \(0.9792~\mathrm{GHz}\), while the target-frequency reflection coefficients degraded to \(-7.839072~\mathrm{dB}\) and \(-5.903831~\mathrm{dB}\), respectively. At \(0.1~\mathrm{mm}\), the documented model produced a dominant minimum at \(1.3976~\mathrm{GHz}\) and \(S_{11}(1.120~\mathrm{GHz})=-6.951098~\mathrm{dB}\). Because the non-monotonic frequency displacement was not verified through mesh or boundary convergence, it is reported as a fixed-model observation rather than a validated physical trend.

At the design frequency, gain decreased from \(0.7515~\mathrm{dBi}\) at \(10~\mathrm{mm}\) to \(0.3534~\mathrm{dBi}\) at \(5~\mathrm{mm}\), and to \(-5.261~\mathrm{dBi}\) at \(0.1~\mathrm{mm}\). Shortening the dipole to \(106.6~\mathrm{mm}\) for the \(5~\mathrm{mm}\) case improved \(S_{11}(1.120~\mathrm{GHz})\) from \(-5.903831~\mathrm{dB}\) to \(-8.000936~\mathrm{dB}\), but did not recover the common \(-10~\mathrm{dB}\) criterion. The study shows that wearable antenna design requires simultaneous control of spacing, detuning, mismatch, tissue absorption, efficiency, and gain.
